\cotefile{34}
\scanpage{3}
\mspage{1}
\uline{Exposé I}\quad \deleted{\uline{Action du groupe de monodromie}} : \uline{Généralités}, \inserted{le} \uline{th. de monodromie par voie arithmétique}.

0. \uline{Introduction}

0.1. Notations \(S=\operatorname{Spec}\deleted{\unsure{R|A}}V\) un trait, \(s\) pt fermé, \(\eta\) point générique
\[
\begin{aligned}
\mathfrak m&=\pi V, & k&=k(s), & K&=k(\eta); & p&=\text{exp. car. de }k(s).\\[-2pt]
&\phantom{=}\shortparallel &&\phantom{=}\shortparallel\\[-2pt]
&\phantom{=}tV &&\phantom{=}V/\mathfrak m
\end{aligned}
\]
Le plus souvent, on suppose \(\deleted{\unsure{R|A}}V\) hensélien, et \(\overline{\mathrm m}\) strictement local i.e. \(k\) sép. clos ;

\deleted{\unsure{\(K=\)|Soit}} \(\overline K\) une clôture séparable de \(K\), \(\bar\eta=\operatorname{Spec}\overline K\) ;
\[
\pi=\operatorname{Gal}(\overline K/K).
\]

\(\ell\) nombre premier. \deleted{Dans la plus grande partie} \deleted{Sauf mention du contraire} \inserted{Pour tout ce \unsure{qu'on fera|qui se fera} de sérieux,} \deleted{\unsure{voir}} on suppose \deleted{L'objet \unsure{du|de}}
\[
\ell\ne\operatorname{car}k
\]

0.2. L'objet du Séminaire est l'étude de certaines représentations \(\ell\)-adiques de \(\pi\), provenant de diverses situations géométriques. Une repr. \(\ell\)-adique de \(\pi\) est un \deleted{\unsure{mor. de}} hom. continu\unsure{s}
\[
\pi\longrightarrow\operatorname{Aut}_{\mathbb Z_\ell}(E),\tag*{(*)}
\]
où \(E\) est un module de t.f. sur \(\mathbb Z_\ell\). On pourra aussi regarder
\[
\pi\longrightarrow\operatorname{Aut}_{\mathbb Q_\ell}(E)
\]
\(E\) module de t.f. sur \(\mathbb Q_\ell\).

\uline{Exemple} 0.3. \(X_K\) schéma algébrique sur \(K\), \inserted{\(i\) entier.} On suppose que \(H^i(X_{\overline K},\mathbb Z/\ell\mathbb Z)\) est un groupe fini, de sorte que
\[
E=H^i(X_{\overline K},\mathbb Z_\ell)\overset{\mathrm{dfn}}{=}\varprojlim_{\nu}H^i(X_{\overline K},\mathbb Z/\ell^{\nu}\mathbb Z)
\]
est un module de t.f. sur \(\mathbb Z_\ell\) ; \unsure{Cette|cette} hyp. \deleted{\illegible{1}} est \deleted{\unsure{K}} vérifiée dans les trois cas suivants :

a) On dispose de résolution des sing. \inserted{sous forme de Hironaka (  )} des schémas alg. de dim \(\leq\dim X\) sur \(\overline K\), -- p. ex. \(K\) de car. nulle

b) \(X\) propre sur \(K\).

c) \(X\) lisse sur \(K\).
% archive sheet number: 2

\scanpage{4}
\mspage{2}
\unsure{Comme} \(\pi\) opère sur \deleted{\(\unsure{H^i(X_{\overline K})}\)} \(E\) par ``transport de structure'', on trouve une représentation \(\ell\)-adique (*)

\uline{Exemple} \unsure{0|2}.4. Soit \(A\) un \deleted{\unsure{schéma en|schéma}} groupe algébrique \inserted{comm.} sur \(K\). \deleted{\unsure{comm.}} alors on regarde
\[
E=T_\ell(A(\overline K))\overset{\mathrm{dfn}}{=}\varprojlim_{\nu}{}_{\ell^\nu}A(\overline K)
\]
\deleted{\unsure{lorsque|Le module}} le groupe ne change pas si on remplace \(A\) par \(\deleted{\unsure{(}}A^\circ_{\mathrm{red}}\) (NB. \(T_\ell\) est un foncteur exact à gauche nul sur les groupes finis), c'est un \deleted{groupe} \uline{module libre sur} \(\mathbb Z_\ell\) de rang égal : \(2\alpha+\deleted{\unsure{2}}\mu\), où \(\alpha\) (\(\mu\)) est la dim. de la ``partie ``abélienne'' resp. \inserted{``partie} ``torique'' de \((A_K)^\circ_{\mathrm{red}}\). Ici encore, \(\pi\) opère sur \(E\) par transport de structure.

\uline{Remarque} 0.5. On peut regarder l'exemple 0.\unsure{3|4} comme étant essentiellement \deleted{un} \inserted{le} cas particulier de l'exemple 0.4, \inserted{pour \(i=1\).} De façon précise, \inserted{on montre qu'}on a un isom canonique
\[
H^1(A_{\overline K},\mathbb Z_\ell)\simeq\deleted{\unsure{H^1(}}\operatorname{Hom}_{\mathbb Z_\ell}(T_\ell(A(\overline K)),\mathbb Z_\ell),
\]
ou ce qui revient au \(\overline{\mathrm m}\), puisque \(T_\ell(A_{\overline K})\) est libre de type fini sur \(\mathbb Z_\ell\) :
\[
T_\ell(A(\overline K))=\operatorname{Hom}_{\mathbb Z_\ell}(H^1(A_{\overline K},\mathbb Z_\ell),\mathbb Z_\ell).
\]

D'autre part, \deleted{dans 0.3, il} si \deleted{\(i=1\)} dans \unsure{0|3}.3. \(X\) est propre et \(i=1\), on trouve
\[
\begin{aligned}
H^1(X_{\overline K},\mathbb Z_\ell(1))&\simeq\varprojlim_\nu H^1(X_{\overline K},\mu_{\ell^\nu})\simeq T_\ell(\operatorname{Pic}(X_{\overline K}))\\
&\simeq T_\ell(\underline{\operatorname{Pic}}_{X/K}(\overline K))
\end{aligned}
\]
% archive sheet number: 3

\scanpage{5}
\mspage{3}
ou si on préfère
\[
H^1(X_{\overline K},\mathbb Z_\ell)\simeq T_\ell(\underline{\operatorname{Pic}}_{X/K}(\overline K))\underset{\mathbb Z_\ell}{\otimes}\mathbb Z_\ell(-1).
\]

0.6. Lorsque \(A\) est un schéma abélien sur \(K\), nous verrons que la connaissance du \(\pi\)-module \(T_\ell(A(\overline K))\) donne des informations \unsure{assez} précises sur la structure du \uline{modèle de Néron} de \(A\) [\quad], -- et inversement la \inserted{considération du mod. de Néron permet d'obtenir des renseignements complémentaires intéressants sur \(T_\ell(A(\overline K))\), -- et partant sur tous les \(H^1(X_{\overline K})\unsure{\bar s|;|f}\), comme \(\pi\)-modules.}
% slanted: continuation of the insertion in the left margin

D'autre part, l'étude des \(\pi\)-modules plus généraux \(H^i(X_{\overline K},\mathbb Z_\ell)\) s'introduit quand on \deleted{\unsure{désigne de}} dispose d'un \deleted{modèle} \inserted{schéma} \(X\) \deleted{\unsure{de}} sur \(S\) de fibre générique \(X_K=X_\eta\), et qu'on désire \deleted{\unsure{étudier}} les relations entre \inserted{\(H^i(X)\),} \(H^i(X_{\unsure{s|\bar s}})\) et \(H^i(X_{\bar\eta})\) ; (p. ex. dans certains cas on pourra montrer \(H^i(X_s)\simeq H^i(X_{\bar\eta})^\pi\)). Cette \deleted{Deligne, utilisant l'étude} étude locale s'avère comme un outil indispensable pour l'étude \inserted{cohomologique} \uline{globale} d'un morphisme
\[
f:X\mathrel{\unsure{-|\longrightarrow}}S
\]
avec \(S\) régulier de dim 1. \deleted{\unsure{Prenant les|Poussant les|Prenant le|Pour}} Dans le cas où \(S=\mathbb P^1_k\), et où \(\deleted{X}f\) ``est'' \deleted{\unsure{déduit|décrit}} un pinceau \inserted{\unsure{linéaires|linéaire}} de sections hyperplanes d'une variété projective lisse \(X_0\) définie sur \(k\), on est dans la situation typique de la \uline{théorie de Lefschetz} [\quad]. S'inspirant d'arguments de Lefschetz, Deligne montrera plus tard que si \(X_0\) est une surface de degré \(\geq4\)
% archive sheet number: 4

\scanpage{6}
\mspage{4}
dans \(\mathbb P^3_k\), alors on a
\[
\mathbb Z\simeq\operatorname{Pic}(\mathbb P^3)\xrightarrow{\sim}\operatorname{Pic}(X_0),
\]
\(k\) étant un corps de car. quelconque.

2. \uline{Rappels sur la structure de} \(\pi\).

On suppose \(S\) hensélien, i.e. \(\deleted{\unsure{R|A}}V\) hensélien.

2.1. Il y a une sous-extension \(\widetilde K_{\deleted{\unsure{e|nr}}}\) de \(\overline K/K\), qui est la plus grande ss-extension non ramifiée (i.e. étale) \deleted{\unsure{dans}} sur \(V\). Si \deleted{\(\unsure{\overline K_e}\) est} \(\overline V\) est la clôture normale de \(V\) ds \(\overline K\), \(\widetilde V=\deleted{\unsure{\overline V_e}=}\overline V\cap\widetilde K_{\deleted{e}}\) la clôture normale dans \(\widetilde K_{\deleted{e}}\), alors \(\widetilde V_{\deleted{\unsure{e}}}\) est un anneau de valuation discrète tel que
\[
\begin{aligned}
&\deleted{\unsure{k|M|\mathfrak m}(\widetilde V)=\mathfrak m\widetilde V,}\\
&\widetilde V\otimes_V k=\bar k
\end{aligned}
\]
est une \deleted{extension \unsure{al}} clôture séparable de \(k\) ; En d'autres termes, \(\widetilde V\) est une \uline{hensélisation stricte} de \(V\) p.r. à \deleted{\unsure{une}} la \deleted{\unsure{co}} clôture séparable \(\bar k\) de \(k\), et \(\widetilde K_{\deleted{e}}\) est le corps des fractions de \(\widetilde V\). On posera \inserted{\(\left\{\begin{aligned}\widetilde S&=\operatorname{Spec}\widetilde V=\text{hensélisé strict de }S\\ \bar s&=\operatorname{Spec}\bar k.\end{aligned}\right.\)}

\margin{\(\circ\mathrel{-}\)} 2.3. On se rappellera que \deleted{\unsure{de sorte|de même}} la donnée de la clôture séparable \(\overline K\) de \(K\) définit canoniquement une clôture séparable \(\bar k\) de \(k\). \deleted{\unsure{Cett|Cette}} Si \(S'\) est un

\deleted{Pour tout \deleted{\unsure{entier le}} entier \(n\) \deleted{premier à} inversible dans \(k\), on a
\[
{}_{n}\bar k^{*}\xleftarrow[\deleted{\unsure{p}}]{\sim}{}_{n}\widetilde V^{*}\xrightarrow{\sim}\overline K^{*}\tag*{2.1.1.}
\]
d'où un isomorphisme canonique
\[
\mu_n(\overline K)\simeq\mu_n(\bar k).
\]
Ces isom sont compatibles avec les morphismes de}
% archive sheet number: 5

\scanpage{7}
\mspage{5}
schéma fini étale sur \(S\), alors on a des isomorphismes
\[
\operatorname{Hom}_{S}(\bar s,S')\xleftarrow{\sim}\operatorname{Hom}_S(\deleted{\operatorname{Spec}}\widetilde S,S')\xrightarrow{\sim}\operatorname{Hom}_S(\deleted{\operatorname{Spec}}\bar\eta,S')
\]
d'où un isom canonique
\[
\operatorname{Hom}_S(\bar s,S')\simeq\operatorname{Hom}_S(\bar\eta,S'),\tag*{(2.3.1)}
\]
fonctoriel en \(S'\) (isom. des \inserted{2} ``foncteurs fibres'' sur la cat. des \unsure{ouv.|rev.|var.} étales de \(S\), définis resp. par les pts géom. \(\unsure{\bar t|\bar\eta|\bar s}\) et \(\bar s\)). Faisant \(S'=\mu_n\), \(n\) \inserted{entier} inversible ds \(k\), on trouve en particulier un isom
\[
\mu_n(\bar s)={}_n\bar k^*\simeq\mu_n(\overline K)={}_n\overline K^*,\tag*{(2.\unsure{3|1}.2)}
\]
d'où par passage à la limite sur les puissances un nb premier \(\ell\), \(\ell\ne\operatorname{car.}k\) :
\[
T_\ell(\bar k^*)\simeq T_\ell(\overline K^*).\tag*{(2.3.3)}
\]
\deleted{On \unsure{notera|note}.} Les deux membres, \deleted{qui} sont des modules libres de rang 1 sur \(\mathbb Z_\ell\), \deleted{\unsure{les ident}} \inserted{qu'on identifie} désormais et désignera par \(\mathbb Z_\ell(1)\). \deleted{\unsure{Ses puissances}}
\[
\mathbb Z_\ell(1)\simeq\ldots\ldots\tag*{(2.3.3)}
\]
\inserted{Sa puiss.} \unsure{tensorielle|tensorielles} \(n\).ième (\(n\in\mathbb Z\)) est notée \(\mathbb Z_\ell(n)\). Donc on a
\[
\mathbb Z_\ell(n)\simeq\varprojlim_\nu\mu_{\ell^\nu}(\bar k^*)^{\otimes n}\simeq\varprojlim_\nu\mu_{\ell^\nu}(\overline K^*)^{\otimes n}.\tag*{(2.3.4)}
\]
Ces isomorphismes \deleted{\unsure{sont compatibles avec}} transforment les opérations de \(\pi\) sur \(\mu_n(\overline K^*)\) \inserted{etc} en celles de \(\pi_0\) sur \(\mu_n(\bar k^*)\) etc\ldots, compte tenu de l'épimorphisme canonique \(\pi\mathrel{\unsure{-|\to}}\pi_0\).
% archive sheet number: 6

\scanpage{8}
\mspage{6}
\margin{\(\circ\!\!-\) \deleted{\unsure{2|3}}}
2.2. La sous-extension \(\overline K_e\) de \(\overline K/K\) est galoisienne, et son groupe de Galois est canoniquement isomorphe au groupe de Galois
\[
\pi_0=\operatorname{Gal}(\bar k/k).\tag*{(2.2.1)}
\]
On obtient ainsi une suite exacte
\[
1\mathrel{\unsure{-|\to}} I\longrightarrow\pi\mathrel{-}\pi_0\mathrel{-}1,\tag*{(2.2.2)}
\]
où
\[
I=\operatorname{Gal}(\overline K/\widetilde K)
\]
est ce qu'on appellé le groupe d'\uline{inertie}. \deleted{\unsure{On}} \deleted{C'est aussi} Si on se rappelle que \(\widetilde K=\) corps des fonctions de \(\widetilde V\), \inserted{et que} \(\overline K\) est une clôture séparable de \(\widetilde K\), on voit donc que \(I\) est le groupe de Galois qui correspond à la situation \deleted{d'un} de l'anneau de valuation discrète \uline{strictement local} \(\widetilde V\). Si \(k\) est sép. clos, on aura \(K=\widetilde K\), \(V=\widetilde V\), \(I=\pi\).

\uline{L'objet du séminaire est} plutôt \uline{l'étude des représentations \(\ell\)-adiques du groupe d'inertie} \(I\). (Ce groupe, ou \unsure{son} image par la représentation \(\ell\)-adique envisagée, \deleted{\unsure{prend alors le|peut alors le|provient de la}} est parfois appelé \uline{groupe de monodromie locale}). Il revient au \(\overline{\mathrm m}\) de dire qu'on s'attache surtout à étudier les repr. \(\ell\)-adiques d'un groupe \(\pi\), lorsque \(k\) est sép. clos. Dans certaines situations cependant, on aura une représentation d'une extension
% archive sheet number: 7

\scanpage{9}
\mspage{7}
(2.2.2), et on ne se bornera pas \deleted{\unsure{et|à}} à étudier sa restriction à \(I\). \deleted{\unsure{C'est|le}} On verra \deleted{\unsure{parfois|toutefois}} que \inserted{parfois} le fait qu'une représentation de \(I\) se prolonge à \(\pi\) impose à cette représentation des restrictions draconiennes.

2.3.

2.4. Rappelons la structure de \(I\). Soit \(p\) l'exposant caractéristique de \(k\), \deleted{Alors} et soit \(P\) \deleted{le} un \deleted{plus grand} sous-\(p\)-groupe de Silow de \(I\) (\(P=1\) si \(p=1\) i.e. \(k\) de car. nulle). Alors \(P\) \uline{est invariant}, i.e. on a une suite exacte
\[
1\mathrel{\unsure{-|\to}}P\longrightarrow I\mathrel{-} I_t\mathrel{-}1,\tag*{(2.4.1)}
\]
où \(I_t\) \inserted{(partie ``tame'' du groupe d'inertie \(I\))} est un groupe profini premier à \(p\), d'autre part on a un isomorphisme canonique
\[
I_t\simeq\prod_{\ell\ne p}\mathbb Z_\ell(1)\tag*{(2.4.2)}
\]

Ce dévissage est obtenu ainsi. Le groupe \(P\) est le groupe de Galois \deleted{de}
\[
P=\operatorname{Gal}(\overline K/\overline K_t),\tag*{(2.4.1)}
\]
où \(\overline K_t\) est la plus \deleted{\unsure{sous|grande}} grande sous-extension \uline{modérément ramifiée} de \(\overline K/K\) \inserted{(rel. à \(V\))}, ou encore de \(\overline K/\widetilde K\) (rel. à \(\widetilde V\)). Celle-ci s'obtient, en termes
% archive sheet number: 8

\scanpage{10}
\mspage{\unsure{8|9|1}}
d'une uniformisante \deleted{\unsure{de}} \deleted{\(\unsure{\eta|g}\)} \(t\) de \(V\) \inserted{donc de \(\widetilde V\)}, comme la réunion des sous-extensions \deleted{abéliennes} Kummeriennes
\[
K_n=\widetilde K(\sqrt[n]{t})\subset\overline K,\tag*{(2.4.4)}
\]
où \(n\) parcourt les entiers \(>0\) premiers à \(p\). Le groupe de Galois de cette extension \(K_n\) est naturellement isomorphe au groupe
\[
\mu_n(\overline K)={}_n\overline K^*
\]
des racines \(n\)ièmes de l'unité dans \(\overline K\), et par passage à la limite sur \(n\) on trouve
\[
I_t=\operatorname{Gal}((\varinjlim_n K_n)/K)\simeq\varprojlim_n\operatorname{Gal}(K_n/K),\deleted{\unsure{\mathbb Z_\ell(1)}},
\]
i.e (2.4.2).

\uline{Remarque} 2.5. On notera que par définition, le sous-groupe \(P\) de \(I\) est stable par \uline{\unsure{tout}} autom. de \(I\), en particulier par les automorphismes induits par les autom. int. de \(\pi\). Donc \(\pi\) opère par passage au quotient sur \(I/p=I_t\) ; \deleted{On constate aussitôt que} \deleted{cette \unsure{opération|op. équivaut}} ds cette opération, le ss-groupe \(I\) de \(\pi\) opère évidemment trivialement, \(I_t\) étant commutatif ; on en trouve en fait une opération de \(\pi/I\simeq\pi_0\) sur \(I_t\). \deleted{\unsure{Compte t}} On vérifie trivialement que
% archive sheet number: 9

\scanpage{11}
\mspage{9}
cette opération, via l'isomorphisme (2.4.\unsure{1|2}), \deleted{n'est} n'est autre que celle qui provient de \deleted{l'isomorphisme de} \inserted{l'opération naturelle} de \(\pi_0\) sur les facteurs \(\mathbb Z_\ell(1)\), via l'expression (2.3.3) de \(\mathbb Z_\ell(1)\).

Bien entendu, pour tout \(\ell\) \deleted{\unsure{on a}} \inserted{il existe} \unsure{un|des} \unsure{isomorphisme|isomorphismes}
\[
\mathbb Z_\ell(1)\simeq\mathbb Z_\ell,
\]
dont le choix revient au choix d'un syst. cohérent de racines \inserted{primitives} \(\ell^\nu\)-ièmes de 1 dans \(\overline K\) ; Mais on fera attention qu'il n'y a pas de façon naturelle de faire ce choix, et que c'est l'expression intrinsèque (2.4.2) qui seule \deleted{\unsure{peut|qui|permet}} tient compte des opérations naturelles de \(\pi_0\) sur le premier membre.
% archive sheet number: 10

\scanpage{12}
\mspage{10}
3. \uline{Comportement pour \deleted{\unsure{\(V\) var|les}} \(S\) variable}.

Soit
\[
f:S'=\operatorname{Spec}(V')\longrightarrow S=\operatorname{Spec}(V,\tag*{(3.1)}
\]
un morphisme \uline{dominant} de traits, \deleted{Alors} donc \(f(s')=s,\ f(\eta')=\eta\). La donnée d' une \deleted{extension} \inserted{clôture} séparable \deleted{\unsure{ment close}} \(\overline K'\) de \(K'\) en définit une \(\overline K\) de \(K\), à savoir la clôture séparable de \(K'\) dans l'extension \(\overline K'\). Alors

\diagram{\((3.2)\)\quad\begin{tikzcd}[ampersand replacement=\&,column sep=small,row sep=small]
K' \arrow[r,dash]\arrow[d,dash] \& \widetilde K' \arrow[r,dash]\arrow[d,dash] \& \overline K' \arrow[d,dash] \\
K \arrow[r,dash] \& \widetilde K \arrow[r,dash] \& \overline K
\end{tikzcd}}

\(\widetilde K\) est contenu dans \(\widetilde K'\), et \(\widetilde V\) dans \(\widetilde V'\). L'hom d'inclusion \(\widetilde V\mathrel{\unsure{\subset|-|\hookrightarrow}}\widetilde V'\) définit aussi un hom \(\bar k\mathrel{\unsure{-|\to}}\bar k'\), correspondant à \deleted{l'hom} \inserted{un} diagramme commutatif de corps.

\diagram{\((3.3)\)\quad\begin{tikzcd}[ampersand replacement=\&,column sep=small,row sep=small]
\deleted{\unsure{R}}V' \arrow[r,dash]\arrow[d,dash] \& \widetilde V' \arrow[d,dash]\\
\deleted{R}V \arrow[r,dash] \& \widetilde V
\end{tikzcd}\qquad\((3.4)\)\quad\begin{tikzcd}[ampersand replacement=\&,column sep=small,row sep=small]
k' \arrow[r,dash]\arrow[d,dash] \& \bar k' \arrow[d,dash]\\
k \arrow[r,dash] \& \widetilde k
\end{tikzcd}}

On trouve ainsi un diagramme commutatif d'extensions, induit par \(\operatorname{Gal}(\overline K'/K)=\pi'\mathrel{-}\operatorname{Gal}(\overline K/K)=\pi\) :

\diagram{\((3.5)\)\quad\begin{tikzcd}[ampersand replacement=\&,column sep=small,row sep=small]
1\arrow[r,dash] \& I'\arrow[r,dash]\arrow[d] \& \pi'\arrow[r,dash]\arrow[d,dash] \& \pi'_0\arrow[r,dash]\arrow[d,dash] \& 1\\
1\arrow[r,dash] \& I\arrow[r,dash] \& \pi\arrow[r,dash] \& \pi_0\arrow[r,dash] \& 1
\end{tikzcd}}

L'hom \(\pi'_0\mathrel{-}\pi_0\) n'est autre que l'hom
\diagram{\((3.6)\)\quad\begin{tikzcd}[ampersand replacement=\&,column sep=small,row sep=small]
\operatorname{Gal}(\bar k'/k')\arrow[r,dash]\arrow[d,equal] \& \operatorname{Gal}(\bar k/k)\arrow[d,equal]\\
\pi'_0 \& \pi_0
\end{tikzcd}}
déduit de (3.4). L'hom \(I'\mathrel{-}I\) s'insère de \(\overline{\mathrm m}\) ds un diagramme de suites exactes

\diagram{\((3,7)\)\quad\begin{tikzcd}[ampersand replacement=\&,column sep=small,row sep=small]
1\arrow[r] \& P'\arrow[r,dash]\arrow[d] \& I'\arrow[r]\arrow[d] \& I'_t\arrow[r,dash]\arrow[d,dash] \& 1\\
1\arrow[r,dash] \& p\arrow[r,dash] \& I\arrow[r,dash] \& I_t\arrow[r,dash] \& 1
\end{tikzcd}}
% archive sheet number: 11

\scanpage{13}
\mspage{11}
Il convient de déterminer l'hom.
\[
I'_t\simeq\prod_{\ell\ne p}\mathbb Z_\ell(1)\longrightarrow\inserted{I_t\simeq}\prod_{\ell\ne p}\mathbb Z_\ell(1)\tag*{(3.8)}
\]
où on identifie, grâce à l'inclusion \(\overline K\subset\overline K'\), les \deleted{\unsure{deux|les}} groupes \(\mathbb Z_\ell(1)\) définis à l'aide de \(\overline K\) ou de \(\overline K'\) (savoir \(T_\ell(\overline K^*)\) et \(T_\ell(\overline K'{}^*)\)).

Pour ceci, on regarde les sous-corps (pour \((n,p)=1\))
\[
\left\{\begin{aligned}
K_n&=\widetilde K(\unsure{\Theta|\theta}_n)\subset\overline K_t&\text{où}\quad&\Theta_n^n=t,\\
K'_n&=\widetilde K'(\Theta'_n)\subset\overline K'_t&\text{où}\quad&\Theta'_n{}^n=t'
\end{aligned}\right.\tag*{(3.9)}
\]
et on regarde comment
\[
\operatorname{Gal}(K'_n/\widetilde K')\simeq\mu_n(\widetilde K')\xleftarrow{\sim}\mu_n(\widetilde K)
\]
(gpe des racines \(n\)ièmes \(\varepsilon\) de 1 ds \(\widetilde K\)) opère sur \(K_n\). Si \(\varepsilon_{K'}\) désigne l'opération sur \(K'_n\) définie par \(\varepsilon\), donc
\[
\varepsilon_{K'_n}(\Theta'_n)=\varepsilon\Theta'_n
\]
on trouve, si
\[
t=\deleted{\unsure{t'}}u't'{}^e\quad\inserted{(e\in\mathbb N,\ u'\in V^*)}\quad
\text{d'où }\deleted{\unsure{=}}\ \boxed{\begin{aligned}&\deleted{\widetilde{\mathfrak m}\widetilde V=\widetilde{\mathfrak m}'\unsure{V'}}\\&\mathfrak m V=\mathfrak m'{}^e\end{aligned}}\tag*{(3.10)}
\]
que
\[
\Theta_n^n=u'(\Theta'_n{}^e)^n
\]
donc, si on pose
\[
\Theta_n=\lambda\Theta'_n{}^e,\quad\inserted{(\lambda\in K'_n)}\tag*{(3.11)}
\]
on aura
\[
\lambda^n=u',\tag*{\deleted{\unsure{(3.12)}}}
\]
donc \deleted{\unsure{reste|d'où}} \(\lambda\in\widetilde V'\) \deleted{\unsure{et}} donc \(\varepsilon_{K'}(\lambda)=\lambda\) et on trouve
\[
\varepsilon_{K'_n}(\Theta_n)=\underset{\substack{\shortparallel\\\lambda}}{\varepsilon_{K'_n}(\lambda)}\varepsilon_{K'_n}(\Theta'_n{}^{\deleted{\unsure{e}}})^e=\lambda(\varepsilon\Theta'_n)^e\overset{\inserted{\text{par (3.10)}}}{=}\varepsilon^e\Theta_n
\]
% archive sheet number: 12

\scanpage{14}
\mspage{12}
en d'autres termes
\[
 \varepsilon_{K'_n}(\Theta_n)=(\varepsilon^e)^{\unsure{n|\cdot}}_{K_n}(\Theta_n)
\]
et par suite
\[
 \varepsilon_{K'_n}\mid K_n=(\varepsilon^e)_{K_n}.
\]

Passant à la limite sur \(n\), on trouve la

\uline{Proposition} 3.12. Si \(\mathfrak mV'=\mathfrak m'{}^e\), alors l'hom. \inserted{\unsure{can|can.}} (3.8) provenant de (3.5) s'identifie à la multiplication par \(e\).

En particulier, si \(e=1\), cet hom. est l'identité, \deleted{\unsure{un|et}} et à fortiori un isomorphisme. \inserted{Ceci s'applique en particulier si \(V'=\widehat V\).} En fait, on a mieux :

\uline{Prop.} 3.13. Supposons \(V'=\widehat V\), \deleted{et \(V\) excellent} \deleted{et \(V\) excellent (i.e. \(\widehat K\) séparable sur \(K\))} \deleted{alors} Alors les hom de suites exactes (3.5) et (3.7) sont des isomorphismes.

\deleted{\inserted{\unsure{disons.|déjà vu.}} \(\pi'\unsure{\mathrel{-}|\longrightarrow}\pi\) est \deleted{\unsure{tjrs|toujours}} surjectif, i.e. si \(K_1\) est une extension finie séparable de \(K\), alors \(K_1\otimes_K\widehat K\) est un corps. \deleted{\unsure{Il faut|Il en résulte}} En effet, la clôture normale \(V_1\) de \(V\) dans \(K_1\) est \uline{finie} \deleted{\unsure{sur \(V\)|\(\widehat V\)}} par suite \(V_1\otimes_V\widehat V=\widehat V_1\) est intègre (car \(V_1\) est un anneau \inserted{de val. discrète;} donc \(K_1\otimes\widehat K=\unsure{\widehat V_1|V_1}\otimes\widehat K\) est un corps.}

\deleted{Comme \(\pi'_0\longrightarrow\pi_0\) est évidemment un isom, et aussi \(I'_t\mathrel{-}I_t\) \deleted{\unsure{par l'égalité|par l'équalité}} on est ramené en fait à examiner \(P'\mathrel{-}P\). Ceci \inserted{(ou directement \(\pi'\mathrel{-}\pi\), \deleted{\unsure{\(I\)|et}} donc \(I'\mathrel{-}I\))} est tjrs injectif, d'après le résultat connu de \deleted{\unsure{approximation|approximations}} \inserted{continuité} des solutions d'une équation algébrique p.r. aux coefficients, qui implique que des équations séparables \uline{voisines} définissent des extensions isomorphes de \(\widehat K\). \margin{\(\widetilde V\subset\widetilde{\widehat V}\subset\widehat{\widetilde V}\)}
% slanted
\deleted{D'autre part, c'est surjectif} \inserted{Comme \(I'\longrightarrow I\) est surjectif} (d'après le résultat \deleted{\unsure{que|qui} \illegible{1} fait l'hypothèse d'excellence}}

\deleted{\uline{Remarque} 3.14. La prop. 3.13. montre que dans la plupart des questions \deleted{\unsure{concernant|correspondant}} du séminaire, on pourrait supposer \(V\) \uline{complet}.}

% archive sheet number: 13

\scanpage{15}
\mspage{13}
\uline{Dém} a) Comme \(\pi'_0\unsure{\longrightarrow|\mathrel{-}}\pi_0\) est évidemment un isom (car \(k'/k\) est l'ext. triviale) il suffit de prouver que \(I'\xrightarrow{\sim}I\), et comme \(I'_t\xrightarrow{\sim}I_t\) (car \(e=1\)) il suffit de prouver que \(P'\xrightarrow{\sim}P\).

b) On sait que \(\pi'\mathrel{-}\pi\) est surjectif, i.e pour tte ext. finie sép. \(K_1/K\), \(K_1\otimes_K\widehat K\) est un \uline{corps} : en effet, la clôture normale \(V_1\) de \(V\) dans \(K_1\) est \uline{finie}, donc \(V_1\otimes_V\widehat V\simeq\widehat V_1\), \deleted{\unsure{est|et}} \inserted{c'est donc} un anneau de val. discrète (\(V_1\) en étant un), donc \(K_1\otimes_K\widehat K=\widehat V_1\otimes_V\widehat K\) est un corps.

\deleted{\unsure{D'autre|Du reste|Dans} Quitte à remplacer \(V,V'\) par \(\widetilde V,\widetilde V'\) et à comparer aussi leurs complétés \illegible{2} \(\widehat{\widetilde V}\simeq\widehat{\widetilde V'}\) de ces derniers, \unsure{on est ramené|on est amené}}

c) On \deleted{\unsure{en|est}} conclut \inserted{de a) et b)}, sur les \inserted{\unsure{hom de|hom. de}} suites exactes, que \(I'\mathrel{-}I\) est surjectif, puis \(P'\mathrel{-}P\) \uline{surjectif}. Il reste à prouver que \(P'\mathrel{-}P\) est injectif, et pour ceci que \(\pi'\mathrel{-}\pi\) est injectif, \deleted{Or} i.e. que toute extension séparable finie \(K'_1\) de \(K'=\widehat K\) provient d'une ext. sép. finie de \(K\). Or ceci résulte du fait connu (``continuité des racines p.r. aux coefficients'') que deux équations séparables \deleted{\unsure{de|dans}} à coefficients dans \(\widehat K\) \uline{voisines} définissent des \deleted{\unsure{corps|extensions}} \(\widehat K\)-Algèbres isomorphes.

% archive sheet number: 14

\scanpage{16}
\mspage{14}
3.1\deleted{5}4. Soit \(\pi'\) un sous-groupe d'indice fini de \(\pi\). Il correspond donc à une extension séparable finie \(K'\) de \(K\), d'où un normalisé \(V'\) de \(V\) dans \(K'\), qui est un anneau de val. discrète hensélien. Si \(k'\) est son corps résiduel, \(k'/k\) est une extension finie, donc \(\pi'_0\unsure{\mathrel{-}|\longrightarrow}\pi'\) \deleted{\unsure{est|a}} \inserted{induit} un \deleted{\unsure{isom|isomes}} de \(\pi'_0\)\inserted{\(=\mathrm{Gal}(\bar k'/k')\)} sur un sous groupe ouvert de \(\pi_0=\mathrm{Gal}(\bar k/k)\). \deleted{D'autre part} \inserted{Il s'ensuit en particulier que \(\pi'_0\mathrel{-}\pi'\) est surjectif, i.e. que} \deleted{on voit que} \inserted{l'on a}
\[
 I'=I\cap\pi' \tag*{(3.1\deleted{5}\unsure{4|5}.1)}
\]

Notons que d'après le théor. général appliqué à \(V'\), on aura pour \unsure{tt|tout} \inserted{nb premier} \(\ell\ne p\)
\[
 [I'/(I',I')]^{\deleted{\unsure{(\ell)|\ell}}}(\ell)\simeq\mathbb Z_\ell(1) \tag*{(3.\unsure{14|15}.2)}
\]
(isom canonique). \deleted{D'autre part}, l'hom \(I'_t\mathrel{-}I_t\) s'identifie, sur \(I_t\), à la multiplication par \(e\) \inserted{comme dessus}, \deleted{qui} lorsque \(k'/k\) est séparable, \deleted{p.ex. que} \inserted{i.e.} l'extension résiduelle de \(\widetilde V'/\widetilde V\) \deleted{séparable} \inserted{triviale}, \deleted{\unsure{donc|dans|Alors}} p. ex. si \(k\) est parfait, alors on peut interpréter \(e\) en termes purement ``\unsure{gruppentheoretische|gruppentheoretischer}'' comme
\[
 e=[I:I']=\mathrm{card}\ I/I'.
\]
Sinon, il faut diviser par le facteur inséparable \(p^s\) de \(k'/k\).

\deleted{3.16. Il \unsure{reste à noter autrement|est intéressant}, lorsque \unsure{le corps|le fait} \(\unsure{I'|\pi'}\) invariant dans \(\pi\) \unsure{est invariant|est quotient} de \(I\), \(\pi'/\pi\) opère sur}

3.15. Lorsque \(\pi'\) est invariant dans \(\pi\) \inserted{(i.e. \(K'/K\) galoisienne)}, il en est de \(\overline{\mathrm m}\) de \deleted{\(\pi'_0\) dans \(\pi_0\), puisque (on ne sait pas \(k'/k\) est une ext. \unsure{quasi-|pure-}galoisienne), donc \unsure{l'hom|l'}} \(I'=I\cap\pi'\) \inserted{dans \(\pi\)} (puisque \(I\) est invariant), \deleted{de \(\unsure{\pi/\pi'|\pi'/\pi}\)} \deleted{\unsure{invariant|invariance}} \(\pi\unsure{?|\ne}\) opère sur \(I'/(I',I')^\wedge\) et par suite sur

% archive sheet number: 15

\scanpage{17}
\mspage{15}
\[
 I'_t=\prod_{\ell\ne p}\bigl(I'/[I',I']\bigr)(\ell)\simeq\prod_{\ell\ne p}\mathbb Z_\ell(1). \tag*{(3.\unsure{15|16}.1)}
\]

\deleted{Ceci posé, \unsure{on|nous} constate aussitôt que les opérations de \(\pi\)} \deleted{Bien entendu, \(I'\) opère trivialement \inserted{sur \deleted{\unsure{\(I\)|\(I'\)}} \(I'_{\mathrm{ab}}\)}, donc c'est le groupe \(\pi/I'\) qui opère sur \(I'_{\unsure{\mathrm{ab}|\mathrm{tab}}}\) et à fortiori sur \(I'_t\). Ceci posé, on constate que \(I/I'\) opère trivialement sur (3.16.1), et que les opérations de \(\pi/I\simeq\pi_0\) sur \(I_t\) ainsi}

\deleted{\unsure{Or|Et}} l'isomorphisme \deleted{de \(\pi\)} (3.16.1) respecte les opérations de \(\pi\) sur les deux mbs, où l'opération sur le deuxième mb. \deleted{\unsure{est|se}} se fait via \(\pi_0\) comme expliqué déjà plus haut.

Ceci ne fait qu'exprimer le fait que l'isom canonique (2\,15.1, est compatible avec ``transport de structure''.

En particulier, on trouve que \(I\deleted{/\unsure{I'|I}}\) opère trivialement sur \(I'_t\), en d'autres termes

\uline{Corollaire} 3.1\deleted{\unsure{7|9}}6. Pour tout ss-groupe \inserted{ouvert} invariant \(I'\) \deleted{d'indice \unsure{fini|fin.}} de \(I\), l'extension naturelle de \(I/I'\) par le quotient \(I'_t\simeq\prod_{\ell\ne p}\mathbb Z_\ell(1)\) de \(I'\) est une extension \uline{centrale}.

% archive sheet number: 16
